\section{Normal local domains in every residue characteristic}

Generation can fail over a two-dimensional normal local domain that
is flat over the integers.  The obstruction is an invertible module
of order two on the characteristic-zero fibre, detected as the
invariant part of an induced representation.

\begin{theorem}\label{thm:normal-domain-counterexample}
\label{thm:normal-surface-counterexample}
For every prime $p$ there is a two-dimensional noetherian normal local
domain $A$ of characteristic zero and residue characteristic $p$ such
that the trivial symmetric-group classes do not generate its induction
ring.  The ring $A$ is flat over $\Z_{(p)}$, and its only square roots
of $1$ are $1$ and $-1$.
\end{theorem}

\begin{proof}
Choose a primitive $p$th root of unity $\zeta$, and put
\[
 O=\Z_{(p)}[\zeta],\qquad \pi=\zeta-1.
\]
The ring $O$ is a discrete valuation ring with uniformizer $\pi$ and
residue field $k=\F_p$: the polynomial $\Phi_p(1+X)$ is Eisenstein at
$p$, and $p$ is a unit times $\pi^{p-1}$.
Choose $u\in O^\times$ and $v\in O$ such that
\[
 z^2+\overline v z+\overline u\in k[z]
\]
is irreducible.  For odd $p$, take $v=0$ and $-\overline u$ a nonsquare;
for $p=2$, take $u=v=1$.  Set
\begin{equation}\label{eq:normal-surface-ring}
 A=\left(O[x,b]/(x(\pi-x)-b(vx+ub))\right)_{(\pi,x,b)}.
\end{equation}
We simultaneously verify the complete construction
\begin{equation}\label{eq:complete-normal-surface-ring}
 \widehat A=\widehat O[[x,b]]/(x(\pi-x)-b(vx+ub)).
\end{equation}
In the argument below, $A$ denotes either ring and $O$ is replaced by
$\widehat O$ in the complete case.

Write $f=x(\pi-x)-b(vx+ub)$.  In either regular local ambient ring,
the associated graded quotient is
\[
 k[\Pi,X,B]/(X\Pi-X^2-\overline v XB-\overline u B^2).
\]
The displayed polynomial is primitive and linear in $\Pi$, because
$X$ is relatively prime to $X^2+\overline v XB+\overline u B^2$.
It is therefore irreducible.  The associated graded ring is a domain,
so separatedness of the maximal-ideal filtration proves that $A$ is
a domain.  It is a Cohen--Macaulay hypersurface of dimension two and
embedding dimension three.
Its quotient by $\pi$ is also a domain: the binary quadratic
$X^2+\overline v XB+\overline u B^2$ is irreducible by the choice of
$u,v$, and the associated graded argument also treats the formal
power-series quotient.  Thus $(\pi)$ is a single height-one prime,
whose localization is regular because its maximal ideal is principal.

Put $C=A[1/\pi]$.  This ring is regular.  Indeed,
\[
 f_x=\pi-2x-vb,\qquad f_b=-vx-2ub,
 \qquad xf_x+bf_b=-\pi x\pmod f.
\]
If a prime not containing $\pi$ contained both partial derivatives,
it would contain $x$, then $b$ by $f=0$ and $u$ a unit, and then
$\pi$ by $f_x$, a contradiction.  Hence one partial derivative is a
unit at every such localization.  This proves regularity of the
hypersurface localization: a derivation sends the square of a prime
into that prime, so $f$ does not belong to the square of the localized
maximal ideal in the regular ambient ring.  The same argument uses
formal partial derivatives for~\eqref{eq:complete-normal-surface-ring}.
All height-one localizations of $A$ are therefore regular, and Serre's
criterion proves normality \citep[Lemma 10.157.4, Tag 031S]{Stacks2026}.  Evaluation at $x=b=0$ gives a retraction
onto $O$, so $A$ has characteristic zero.  Its torsion-freeness over
$\Z_{(p)}$ gives flatness.  In a characteristic-zero domain,
$(w-1)(w+1)=0$ forces $w=1$ or $w=-1$.

Consider $I=(x,b)$.  This is a height-one prime, since $A/I=O$.
It is nonprincipal: $x,b$ are independent in the cotangent space
$\mathfrak m/\mathfrak m^2$, and hence in $I/\mathfrak m I$, so two
generators are necessary.  We claim that $I_C$ is still nonprincipal.
If it had a generator, clearing powers of $\pi$ would give a generator
$r\in A$.  Strip off the largest power of $\pi$ dividing $r$, which
exists by the Krull intersection theorem
\citep[Lemma 10.51.4, Tag 00IP]{Stacks2026}, and write the result as
$r_0\notin\pi A$.  The ideal $(r_0)$ is $\pi$-saturated: if
$\pi s=r_0t$, primality of $(\pi)$ gives $t=\pi t_1$, and cancellation
gives $s=r_0t_1$.  Iteration treats every power of $\pi$.
The prime $I$ is also $\pi$-saturated, because $\pi\notin I$.
Contracting $I_C=r_0C$ to $A$ would therefore give $I=(r_0)$,
a contradiction.

The defining equation gives
\[
 I_C^2=x\bigl(x,b,u^{-1}(\pi-x-vb)\bigr)C=xC.
\]
The last equality uses that $u,\pi$ are units in $C$.  Thus $I_C$
is invertible and its Picard class has exact order two.
Put $J=(b,\pi-x)C$.  Similarly,
\[
 I_CJ=b(b,x,\pi-x,vx+ub)C=bC.
\]
Hence $J$ is a nontrivial invertible module, also of order two.

On $L=A^2$ define
\[
 M=\begin{pmatrix}x&b\\vx+ub&\pi-x\end{pmatrix},
 \qquad T=I+M.
\]
The trace and determinant give $M^2=\pi M$.  Therefore
\[
 T^p
 =I+\left(\sum_{j=1}^{p}\binom pj\pi^{j-1}\right)M=I,
\]
because the coefficient times $\pi$ is $(1+\pi)^p-1=0$ and $O$ is
a domain.  Let a generator of $H=C_p$ act by $T$, embed $H$ in $S_p$
as the subgroup generated by a $p$-cycle, and set
\[
 N=\operatorname{Ind}_{H}^{S_p}L.
\]
This is a finite free $A$-lattice.

Over $C$, the $H$-invariant part of $L$ is the image of the idempotent
\[
 I-\pi^{-1}M
 =\pi^{-1}\begin{pmatrix}\pi-x&-b\\-vx-ub&x\end{pmatrix}.
\]
Projection to its first coordinate identifies this image with $J$.
The image ideal is $(\pi-x,b)C$; injectivity follows because an
invariant vector $(0,z)$ satisfies $bz=0$, with $C$ a domain and
$b\ne0$.  Induction for finite groups is also coinduction, so
\[
 N_C^{S_p}\simeq L_C^H\simeq J.
\]
Every positive integer is invertible in $C$.  The averaging projector
therefore makes invariants an exact functor to finite projective
$C$-modules.  Extension of scalars, invariants, and determinant give
an additive homomorphism
\[
 K_0(\Lat_A(S_p))\longrightarrow K_0(C)
                       \xrightarrow{\det}\operatorname{Pic}(C).
\]
For a Young subgroup $S_\lambda$, the invariants of
$C[S_p/S_\lambda]$ form a free line, since invariant coefficients are
constant on the transitive coset set.  The homomorphism vanishes on
every degree-$p$ monomial in the trivial classes, but its value on
$[N]$ is $[J]\ne0$.  This proves failure of generation for both rings.
\end{proof}

\begin{example}
For $p=2$ one may take $u=v=1$, $\zeta=-1$, and $\pi=-2$.  The example
is
\[
 A=\left(\Z_{(2)}[x,b]/(x(x+2)+xb+b^2)\right)_{(2,x,b)},
 \qquad
 T=\begin{pmatrix}1+x&b\\x+b&-1-x\end{pmatrix}.
\]
Thus a rank-two involution already detects the obstruction for $S_2$.
\end{example}

\begin{corollary}\label{cor:complete-normal-domain-counterexample}
The counterexamples in Theorem~\ref{thm:normal-domain-counterexample}
can be chosen complete and two-dimensional.
\end{corollary}

\begin{proof}
Use~\eqref{eq:complete-normal-surface-ring}.  The preceding proof
verified its domain property, normality, surviving nonprincipal
invertible module, and induced representation together with those of
the polynomial-local example.
\end{proof}

\begin{corollary}\label{cor:minimal-normal-dimension}
Dimension two is the least possible dimension of a noetherian normal
local characteristic-zero counterexample to surjectivity of $\Phi$.
\end{corollary}

\begin{proof}
A zero-dimensional normal local domain is a field, and a
one-dimensional noetherian normal local domain is a discrete valuation
ring \citep[Lemma 10.119.7, Tag 00PD]{Stacks2026}.  In characteristic zero both satisfy the regular-local positive
theorem.  Theorem~\ref{thm:normal-domain-counterexample} supplies
counterexamples in dimension two.
\end{proof}
